% Grado 9 -- generado por tools/build_tex.py a partir de raw/grado_09.txt
\section{Grado 9}\label{grado:9}

% PDF p. 30
\dba{1}{Comprende que el movimiento de un cuerpo, en un marco de referencia inercial dado, se puede describir con gráficos y predecir por medio de expresiones matemáticas.}

\evidencias

\begin{itemize}
  \item Describe el movimiento de un cuerpo (rectilíneo uniforme y uniformemente acelerado, en dos dimensiones – circular uniforme y parabólico) en gráficos que relacionan el desplazamiento, la velocidad y la aceleración en función del tiempo.
  \item Predice el movimiento de un cuerpo a partir de las expresiones matemáticas con las que se relaciona, según el caso, la distancia recorrida, la velocidad y la aceleración en función del tiempo.
  \item Identifica las modificaciones necesarias en la descripción del movimiento de un cuerpo, representada en gráficos, cuando se cambia de marco de referencia.
\end{itemize}

\ejemplo

\fig{g09_p30_1}{Ejemplo (solo imagen): un pasajero dentro de un bus en movimiento deja caer una pelota. Para él, dentro del bus, la pelota cae en línea recta vertical; para un peatón en la acera, la misma pelota describe una trayectoria parabólica. Ilustra que la descripción del movimiento depende del marco de referencia.}


